\section{Conclusion}
\label{sec:conclusion}

This paper compared five MPPT algorithms spanning four genuinely distinct
design paradigms -- perturbative, rule-based, learning-based, and
model-based nonlinear control -- under a common, open-source, statistically
rigorous simulation framework, validated against a real commercial PV
datasheet. Sliding mode control achieved the best steady-state tracking
efficiency and lowest oscillation at near-baseline computational cost. A
reduced $5\times5$ fuzzy rule base statistically matched the full
$7\times7$ base's performance, a genuine rule-reduction contribution given
fuzzy logic's substantial computational overhead. Most significantly,
every algorithm across all four paradigms -- including the model-based SMC
controller with its Lyapunov stability guarantee -- failed to reach the
true global maximum power point in the majority of complex partial-shading
patterns tested, with the failure mode differing meaningfully by paradigm:
the perturbative and model-based algorithms converge to a shared local
maximum, while Q-learning (never trained on multi-module shading) fails in
a qualitatively worse way, moving to and oscillating around a
substantially lower-power operating point instead.

\subsection{Limitations}
Several limitations should inform how these results are read and motivate
future work. First, the boost converter enters discontinuous conduction
mode (DCM) at light load (Section~\ref{sec:converter-model}), where the
small-signal model every algorithm's duty-cycle-perturbation reasoning
implicitly relies on no longer holds; whether this contributes to any
algorithm's degraded performance specifically in the low-irradiance and
partial-shading scenarios has not been isolated from the algorithms'
intrinsic global-search limitations, and would require a dedicated
ablation to disentangle. Second, the two-diode model's temperature
dependence is implicit (through the thermal-voltage term) rather than
calibrated against the datasheet's explicit temperature coefficients
(Section~\ref{sec:pv-model}); the NOCT and low-irradiance curve validation
CLAUDE.md's acceptance criteria specify is not yet complete, though the
underlying datasheet data needed for it is now available. Third,
Q-learning's tabular state discretization and training protocol
(Section~\ref{sec:q-learning-algorithm}) are one specific design point in a
much larger learning-based MPPT design space; whether a different state
representation, function approximation, or a training distribution that
includes partial shading would close the gap observed in
Section~\ref{sec:partial-shading-results} is an open question this paper
does not attempt to answer, since doing so changes the fair-comparison
premise of evaluating the same tabular design used throughout. Fourth, this
is a simulation-only study; while the two-diode model is quantitatively
validated against a real datasheet (Section~\ref{sec:model-validation}),
hardware validation would strengthen any of these conclusions further.

\subsection{Future Work}
The clearest direction suggested by this paper's own results is a deeper,
dedicated study of learning-based MPPT under partial shading -- explicitly
training on shading patterns rather than only flat irradiance, and
comparing tabular Q-learning against function-approximation and deep-RL
alternatives -- since the generalization gap found here is a training-data
problem, not necessarily a fundamental limitation of learning-based control.
A second, complementary direction is a fair, reproducible benchmark of
metaheuristic global-search MPPT methods (which this paper deliberately
excludes, to avoid the saturated "yet another bio-inspired optimizer"
pattern) against this paper's partial-shading scenarios directly, since
metaheuristics are specifically designed for the global-search problem this
paper's results show every non-metaheuristic paradigm struggles with.
Finally, the DCM/light-load confound flagged above is a concrete,
answerable question -- rerunning the low-irradiance and partial-shading
scenarios with a converter design that stays in CCM across the full
operating range would isolate whether DCM sensitivity contributes to any
algorithm's measured performance independent of its global-search
capability.

\subsection*{Code and Data Availability}
% Placeholder -- fill in once the GitHub repository and Zenodo archive are
% public. Do not fill in a DOI or repository URL here until they actually
% exist (RESEARCH_PROGRAM.md's standing rule against unproduced specifics
% applies to this statement as much as to any numeric result).
All simulation code, raw Monte Carlo results, and figures are released
under the MIT License at \texttt{[GitHub repository URL -- fill in]} and
archived with a citable DOI at \texttt{[Zenodo DOI -- fill in once
minted]}.
